\subsection{Research questions and observable answers}
\label{sec:research-questions}

We distinguish eight questions about behavior in an economic game. The first concerns strategic reasoning: whether a player considers competing interests and executes feasible decisions. The second concerns opponent modeling: whether its prospective statements about others have checkable targets and agree with subsequent opportunities or actions. The third concerns cooperation, negotiation and deception. The fourth concerns sustained plans and memory. The fifth concerns adaptation after feedback. The sixth concerns asset allocation, liquidity and risk. The seventh concerns collective patterns produced by interaction. The eighth concerns whether the measurements are valid and generalize beyond the observed configuration.

These questions require different evidence. A settled exchange is observable, but its utility advantage requires a comparator. A legal payment witness settles a financing question, but does not maximize expected victory. A recurring validator error is observable, but does not identify the effect of feedback on learning. We therefore declare the object and endpoint before assigning a behavioral interpretation. Table~\ref{tab:question-contract} links each broad question to the narrower empirical test performed here. Unmeasured components remain explicit rather than being inferred from the amount of descriptive detail.

\begin{table}[t]
\centering\small
\begin{tabularx}{\linewidth}{L{0.21\linewidth}YY}
\toprule
Question & Observable target & Interpretation boundary\\
\midrule
Q1. Strategy & Offer consequences, rival-sensitive responses, and executable financing/construction & Feasibility and observed tradeoffs, not optimal expected policy.\\
Q2. Opponents & Explicit prospective claims in a fixed sample and their identifiable endpoints & Forecasts of actions or resources, not direct access to others' beliefs.\\
Q3. Negotiation & Complete offer/response threads, settlement terms, and contextual speech & Activity and immediate consequences, not a causal persuasion or fairness score.\\
Q4. Plans and memory & Stated plans, later opportunities, execution and recorded information & Observed follow-through, not a randomized memory benefit.\\
Q5. Adaptation & Corrections and repeated errors at later submitted actions & Within-record recurrence, not training or demonstrated durable learning.\\
Q6. Economics & Expenditure, receipts, collateral access, liquidation and exposure & Realized accounting and local constraints, not ex ante asset value or risk-adjusted return.\\
Q7. Collective behavior & Cash/deed interaction networks, concentration and coordination cases & Population flows and documented participation, not a latent collusion rate.\\
Q8. Validity & Replay, source checks, certificate tests and reference-evaluation quality & Validity for stated targets; transfer to new models or opponents needs separate evaluation.\\
\bottomrule
\end{tabularx}
\caption{Question-to-measurement contract. Each row distinguishes the broad research question from the target that this study can evaluate. No row defines a general intelligence score.}
\label{tab:question-contract}
\end{table}

Three broader questions organize the synthesis. First, do strategic explanations correspond to strategic execution? We compare propositions, submitted plans, received actions and consequences separately. Second, does general capability predict multi-agent performance? That association is not identified by four heterogeneous configurations without an independent capability axis. Third, can performance be separated from luck and opponent mistakes? We identify particular errors and transfer mechanisms, but neither a shared seed nor a local legal alternative identifies a counterfactual winner. These are limits on the estimands, not evidence that the broader capabilities are absent.
